\chapter{Appendix F: Computer-Assisted Verification Package}
\label{app:cap_package}

\section{Overview of the CAP Dependency}
This proof utilizes a Computer-Assisted Proof (CAP) specifically for the **Intermediate Regime Bridge** ($SU(2)/SU(3)$ cases). The CAP is designed to verify the "Tube Contraction" condition (Theorem 3.5), ensuring no phase transition occurs between the strong and weak coupling analytical domains.

\textbf{Logical Status:}
\begin{itemize}
    \item \textbf{Analytic Spine:} The strong and weak coupling gaps are proven purely analytically (Appendices B \& C).
    \item \textbf{CAP Bridge:} The absence of singularities in the crossover region region $[ \beta_{strong}, \beta_{weak} ]$ is the specific claim certified by CAP.
\end{itemize}

\section{Verification Artifacts}
The following artifacts are provided in the supplementary material to allow for independent auditing:

\subsection{1. Interval Arithmetic Kernel}
\texttt{lib/interval.hpp}: A C++ header-only library implementing IEEE 754 compliant interval arithmetic with directed rounding modes.
\begin{itemize}
    \item \textbf{Function:} Provides guaranteed inclusion bounds $[x_{lo}, x_{hi}]$ for elementary operations.
    \item \textbf{Certification:} Unit tests cover edge cases (signed zeros, infinity, NaN propagation).
\end{itemize}

\subsection{2. Renormalization Group Map Implementation}
\texttt{src/rg\_step.cpp}: Implements the Kadanoff-Wilson block-spin transformation $R(T)$.
\begin{itemize}
    \item \textbf{Input:} A "Tube" definition (a set of interval bounds on polymer coefficients).
    \item \textbf{Output:} The image of the tube under one coarse-graining step.
    \item \textbf{Check:} Verifies $R(\text{Tube}) \subset \text{Tube}$ (contraction condition).
\end{itemize}

\subsection{3. Execution Scripts & Logs}
\texttt{verify\_gap.sh}: A shell script that compiles the verifier and runs it across the specified $\beta$ range.
\texttt{logs/verification\_chain.log}: A cryptographically hased log of the successful contraction checks for each step along the RG trajectory.

\section{How to Verify}
To audit the proof:
\begin{enumerate}
    \item Inspect \texttt{src/rg\_step.cpp} to confirm it faithfully implements the analytical RG map defined in Eq (3.24).
    \item Run \texttt{./verify\_gap.sh} on a standard Linux environment.
    \item The program generates a "Certificate of Contraction" which implies the absence of critical points (Theorem 3.5).
\end{enumerate}
